% Chapter 6: The Edit-Test-Commit Loop
\chapter{The Edit-Test-Commit Loop}
\label{ch:testing}

\begin{itshape}
Claude just generated 200 lines of code that looks correct.
The syntax is clean, the variable names are reasonable, the logic reads well.
How do you know it actually works?
\end{itshape}

AI-generated code has a specific failure mode that human-written code does not:
it \emph{looks} correct. The appearance of correctness is precisely
why verification is essential---the bugs are subtle, not obvious.

\section{Verification Is King}
\label{sec:verification}

\convergence{``Give Claude a way to verify its work'' is Anthropic's stated
single highest-leverage practice.\sourceurl{https://code.claude.com/docs/en/best-practices}
This aligns exactly with what practitioners discover independently:
the projects that succeed are the ones with verification at every layer.}

\begin{whybox}[Why Verification Is the Highest-Leverage Practice]
Claude Code's creator, Boris Cherny, suggests that verification improves
code quality by 2--3$\times$---a plausible but hard-to-measure claim
that aligns with practitioner experience. The mechanism is straightforward:
when Claude can run tests, check types, and validate output against
concrete criteria, it catches its own mistakes before you see them.
The verification loop turns a single-pass generation into an iterative
refinement process.

Without verification criteria, you are the only quality gate.
With them, Claude becomes a self-correcting system.
\end{whybox}

\subsection{The 6-Layer Validation Architecture}

Layer defenses so that no single failure mode goes undetected:\index{testing!6-layer model}

\begin{enumerate}
  \item \textbf{Type Safety} --- Catch type errors at compile/lint time.
    Python type hints with mypy, pandas-stubs for DataFrame types, pandera for schema validation.

  \item \textbf{Input Validation} --- Check preconditions at function entry.
    Fail fast with explicit errors (\cref{ch:principles}).

  \item \textbf{Unit Tests} --- Test each function in isolation.
    Happy path + error cases + edge cases.
    Target: 80\%+ coverage for modules (phase-appropriate---see \cref{sec:phases}).

  \item \textbf{Integration Tests} --- Test multi-function workflows.
    Verify components work together with realistic data.

  \item \textbf{End-to-End Tests} --- Test complete user workflows
    from input to output. Verify the system meets requirements.

  \item \textbf{Property-Based Tests} --- Test invariants that should always hold.
    Generate random inputs. Catch edge cases you did not think of.
\end{enumerate}

\practitioner{Property-based tests are underused. A single Hypothesis test
can replace dozens of hand-written edge-case tests:}\index{testing!property-based}

\begin{minted}{python}
from hypothesis import given, settings
from hypothesis import strategies as st
import pandas as pd

@given(st.lists(
    st.floats(allow_nan=True, allow_infinity=False),
    min_size=1, max_size=100
))
@settings(max_examples=200)
def test_feature_builder_invariants(values):
    df = pd.DataFrame({"amount": values})
    result = build_features(df)
    # Schema invariant: output columns never change
    assert set(result.columns) == {"amount", "log_amount", "is_missing"}
    # Null invariant: NaN inputs produce is_missing=True
    assert (df["amount"].isna() == result["is_missing"]).all()
    # Range invariant: log_amount is never negative
    assert (result["log_amount"].dropna() >= 0).all()
\end{minted}

\margintip{Layers 1--2 are cheap to add to any project today.
Layers 3--4 are the baseline for production.
Layers 5--6 are for critical systems.}

\subsection{The Missing Layer: Domain Correctness}
\label{subsec:domain-correctness}

\marginnote[Warning]{Claude's most dangerous failure:
code that passes all tests but is
semantically wrong.}

The 6-layer model catches structural errors---type mismatches,
null handling, regression failures. But it does not catch
\emph{domain-correctness} failures: code that is syntactically
valid, passes all tests, and produces the wrong answer.

Examples from practice:

\begin{itemize}
  \item \textbf{Data leakage.} A feature engineering function uses
    future values during training. Tests pass because the function
    is deterministic. The model achieves unrealistic accuracy.
  \item \textbf{Wrong aggregation.} A revenue calculation sums
    instead of averaging across time periods. Tests pass because
    the function produces a number. The number is wrong.
  \item \textbf{Survivorship bias.} A cohort analysis excludes
    records that were deleted. Tests pass because the query runs.
    The results are misleading.
  \item \textbf{Silent unit mismatch.} A function mixes daily and
    monthly rates. Tests pass because both are floats.
    The financial model is off by 30$\times$.
\end{itemize}

\noindent These failures share a pattern: structural tests verify that
the code \emph{runs}, but not that it \emph{means} what you intended.

\subsection*{Domain Invariant Tests}

The fix is a class of tests that encode domain knowledge,
not just interface contracts:

\begin{minted}{python}
def test_no_future_leakage():
    """Features at time t must use only data from t-1 or earlier."""
    df = make_time_series(n=100)
    features = build_features(df, as_of=50)
    # No feature value should depend on rows 51-100
    assert features.index.max() <= 50

def test_revenue_units():
    """Monthly revenue * 12 should approximate annual revenue."""
    monthly = compute_monthly_revenue(test_data)
    annual = compute_annual_revenue(test_data)
    # Within 5% tolerance (seasonal variation)
    assert abs(monthly.sum() * 12 - annual.sum()) / annual.sum() < 0.05

def test_cohort_completeness():
    """Cohort analysis must account for all original records."""
    original_count = len(raw_data)
    cohort_counts = cohort_analysis(raw_data)
    assert cohort_counts.sum() == original_count
\end{minted}

\practitioner{When asking Claude to generate tests, explicitly request
domain invariant tests in addition to structural tests:
``Write tests that verify the business logic is correct, not just that
the code runs. Include: leakage checks, unit consistency,
and completeness invariants.''}

\begin{keyinsight}
The 6-layer model is necessary but not sufficient for domains where
``correct output'' requires judgment. Domain invariant tests encode
\emph{your} knowledge of what the answer should look like. Claude
can generate the test structure, but you must specify the invariants---
because Claude does not know that monthly revenue times twelve should
approximate annual revenue. You do.
\end{keyinsight}

% -------------------------------------------------------------------
\section{Phase-Appropriate Standards}
\label{sec:phases}

\practitioner{Not all code needs the same rigor.
Applying production standards to a prototype kills velocity.
Applying prototype standards to production kills reliability.}\index{testing!phase-appropriate}

\begin{fullwidth}
\begin{tabular}{@{}lp{3cm}p{3cm}p{4cm}@{}}
\toprule
\textbf{Phase} & \textbf{Testing} & \textbf{Code Quality} & \textbf{Transition Criteria} \\
\midrule
Exploration & Manual OK & Long functions OK & Hypothesis validated \\
Development & Unit + integration & Style enforced, type hints & Coverage >60\%, code review \\
Production & Full 6-layer & Strict lint, immutability & Coverage >80\%, zero critical warnings \\
\bottomrule
\end{tabular}
\end{fullwidth}

\subsection{Making Phase Transitions Explicit}

\practitioner{Write phase transitions in CLAUDE.md:}

\begin{minted}{markdown}
## Phase: EXPLORATION (until March 15)
- Tests: manual OK for now
- Coverage target: none
- Graduation criteria:
  - [ ] Model validated on holdout set
  - [ ] No data leakage verified
  - [ ] Feature distributions documented
  - [ ] Dependencies locked
  When all checked: switch to DEVELOPMENT phase
\end{minted}

The key insight: premature rigor kills exploration,
but sloppy production kills credibility.
The solution is not to choose one standard---it is to be explicit about
which standard applies \emph{right now} and when to transition.

% -------------------------------------------------------------------
\section{Test-First with Claude}
\label{sec:test-first}

\marginnote[Official]{``Give Claude a way to verify its own work.''
For greenfield, this means: provide test cases in your initial prompt,
not after the code is written.\sourceurl{https://code.claude.com/docs/en/best-practices}}

\practitioner{The test-first pattern produces higher-quality code
than code-first:}\index{testing!test-first pattern}

\begin{enumerate}
  \item Describe the \textbf{interface}: inputs, outputs, edge cases.
  \item Claude writes \textbf{tests} based on the interface.
  \item Claude writes \textbf{implementation} that passes the tests.
  \item Tests run \textbf{automatically} via hooks.
\end{enumerate}

\begin{minted}{text}
Prompt: "Create a FeatureValidator class:
- Takes a DataFrame and a schema dict
- Validates column types, value ranges, null counts
- Returns ValidationResult with errors list
- Raises ValueError if required columns are missing

Write tests first, then implementation."
\end{minted}

This works because Claude excels at test generation when given
clear specifications. The tests then constrain the implementation,
preventing the ``looks correct but is subtly wrong'' failure mode.

\marginnote[Cross-Ref]{For tests as exploration tools---discovering
requirements, not just verifying them---see \cref{ch:thinking-partner}.}

% -------------------------------------------------------------------
\section{Verification Criteria in CLAUDE.md}
\label{sec:verification-criteria}

\begin{fullwidth}
\begin{beforeafter}
\before
\begin{minted}{markdown}
## Testing
We use pytest. Tests are important.
\end{minted}

\after
\begin{minted}{markdown}
## Verification
- Always run tests after code changes
- Never commit without passing lint + tests
- Include edge case tests for every new function
- Coverage target: 80% for src/ modules
\end{minted}
\end{beforeafter}
\end{fullwidth}

\convergence{These lines in CLAUDE.md improve every interaction
by giving Claude a concrete way to validate its own work.}

% -------------------------------------------------------------------
\section{Quality Gates via Hooks}
\label{sec:quality-gates}

Hooks enforce what CLAUDE.md can only suggest (see \cref{ch:extending}
for the full treatment). Three hooks that prevent quality debt:

\begin{minted}{json}
{
  "hooks": {
    "PostToolUse": [{
      "matcher": "Edit|Write",
      "hooks": [{
        "type": "command",
        "command": ".claude/hooks/lint-on-write.sh",
        "description": "Auto-lint on write"
      }]
    }],
    "PreToolUse": [{
      "matcher": "Bash",
      "hooks": [{
        "type": "command",
        "command": ".claude/hooks/pre-commit-gate.sh",
        "description": "Block bad commits"
      }]
    }]
  }
}
\end{minted}

\begin{whybox}[Why Hooks Beat Instructions]
CLAUDE.md: ``Always run tests before committing.''
Claude may forget, especially in long sessions with degraded context.

Hook: \code{PreToolUse} matching \code{Bash} $\rightarrow$ gate on
\code{git commit} commands, run \code{pytest} first.
Tests run every time, regardless of context state. Cannot be forgotten.

The principle: if a standard is non-negotiable, make it a hook.
If it is advisory, put it in CLAUDE.md.
\end{whybox}

% -------------------------------------------------------------------
\section{Quick Reference}

\begin{itemize}
  \item AI-generated code \emph{looks} correct---verification is essential
  \item 6-layer validation: types $\rightarrow$ input checks $\rightarrow$ unit $\rightarrow$ integration $\rightarrow$ E2E $\rightarrow$ property
  \item Match rigor to phase: exploration $\rightarrow$ development $\rightarrow$ production
  \item Test-first: describe interface $\rightarrow$ tests $\rightarrow$ implementation
  \item Verification criteria in CLAUDE.md improve every session
  \item Hooks enforce non-negotiable standards
\end{itemize}

% -------------------------------------------------------------------
\begin{trybox}
\begin{enumerate}
  \item Add these three lines to your CLAUDE.md:
    ``Always run \code{pytest tests/} after code changes.
    Never commit without passing \code{ruff check src/ \&\& mypy src/}.
    Include edge case tests (empty DataFrames, NaN-heavy inputs)
    for every new function.''
  \item Add a PreToolUse hook (matching \code{Bash}) that gates commits on
    \code{pytest tests/ \&\& mypy src/}.
  \item Track: does Claude's test coverage improve over the next week?
\end{enumerate}
\end{trybox}
